% =====================================================================
%  LISTAS -- Manual, seções 2.1.9 a 2.1.12
% ---------------------------------------------------------------------
%  São OPCIONAIS, mas tornam-se OBRIGATÓRIAS quando houver MAIS DE 5
%  itens do tipo no trabalho. Se tiver 5 ou menos, você pode remover a
%  lista colocando % no início da(s) linha(s) correspondente(s).
%
%  As listas de ilustrações, quadros, tabelas e códigos são geradas
%  AUTOMATICAMENTE a partir dos \caption{} do texto.
%  As listas de siglas e de símbolos são escritas À MÃO, logo abaixo.
% =====================================================================

% --- Lista de ilustrações (figuras, gráficos, fluxogramas, fotos...) ---
\pdfbookmark[0]{\listfigurename}{lof}
\listoffigures*
\cleardoublepage

% --- Lista de quadros (quadro = conteúdo textual, com bordas fechadas) ---
\pdfbookmark[0]{\listofquadrosname}{loq}
\listofquadros*
\cleardoublepage

% --- Lista de tabelas (tabela = dados numéricos, bordas laterais abertas) ---
\pdfbookmark[0]{\listtablename}{lot}
\listoftables*
\cleardoublepage

% --- Lista de códigos (útil no relatório técnico) ---------------------
\ifrelatorio
  \pdfbookmark[0]{\listofcodigosname}{loc}
  \listofcodigos*
  \cleardoublepage
\fi

% --- Lista de abreviaturas e siglas (em ORDEM ALFABÉTICA) -------------
% Na primeira vez que a sigla aparece no texto, escreva o nome completo
% seguido da sigla entre parênteses: "Associação Brasileira de Normas
% Técnicas (ABNT)". Depois, use apenas a sigla.
\begin{siglas}
  \item[ABNT] Associação Brasileira de Normas Técnicas
  \item[CT] Caso de Teste
  \item[IA] Inteligência Artificial
  \item[LGPD] Lei Geral de Proteção de Dados Pessoais
  \item[LLM] \textit{Large Language Model} (modelo de linguagem de grande porte)
  \item[MoSCoW] \textit{Must, Should, Could, Won't} (método de priorização)
  \item[NoSQL] \textit{Not only SQL} (banco de dados não relacional)
  \item[RF] Requisito Funcional
  \item[RN] Regra de Negócio
  \item[RNF] Requisito Não Funcional
  \item[TCC] Trabalho de Conclusão de Curso
  \item[UC] Caso de Uso (\textit{Use Case})
  \item[UML] \textit{Unified Modeling Language}
  \item[WCAG] \textit{Web Content Accessibility Guidelines}
  \item[XP] Pontos de experiência (\textit{experience points})
\end{siglas}

% --- Lista de símbolos (na ORDEM em que aparecem no texto) ------------
% Este trabalho não usa símbolos matemáticos; lista desativada.
% \begin{simbolos}
%   \item[$ \sigma $] Tensão normal
%   \item[$ F $] Força aplicada
%   \item[$ A $] Área da seção transversal
%   \item[$ \in $] Pertence
% \end{simbolos}
